Congressional redistricting in the United States operates under a fundamental
structural constraint: 435 single-member districts elected by winner-take-all
plurality voting.  This combination creates systematic incentives for partisan
manipulation.  A party that controls the redistricting process can engineer
district boundaries to waste opposition votes---packing them into a few
overwhelmingly safe districts, or cracking them across many districts where
they form ineffective minorities \citep{stephanopoulos2015partisan}.  The
result is a persistent gap between a party's statewide vote share and its seat
share, a gap that is structurally guaranteed to favour whichever party drew
the map.

Multi-member districts (MMDs) break this dynamic at the source.  When a
district elects four representatives and seats are allocated proportionally
within the district, a party that wins 40\% of the district's votes receives
approximately 40\% of its seats.  No boundary can be drawn to deny a large
minority party its seats once the allocation rule is changed from
winner-take-all to proportional.  \citet{mcgann2016gerrymandering} argue that
MMDs are structurally gerrymandering-proof: the proportional allocation
mechanism makes packing and cracking strategies geometrically ineffective.

The United States has historical experience with MMDs.  Illinois elected
state legislators from three-member districts using Cumulative Voting from
1870 to 1980 \citep{cole2001cumulative}.  Several U.S. cities (Cambridge,
Massachusetts; Portland, Maine) currently use STV for municipal elections.
The Federal Apportionment Act of 1842 was the first statute to mandate
single-member districts for U.S. House elections; the Constitution itself
does not require single-member districts.

\subsection{This Paper}

We ask a concrete algorithmic question: can the bisect recursive bisection
framework generate multi-member congressional districts, and if so, what are
the compactness and representation consequences?  We answer affirmatively.
Bisect generates MMDs by adjusting the district count $k$:
\begin{itemize}
  \item For California (52 seats), $k = 13$ districts of 4 seats each, or
        $k = 26$ districts of 2 seats each.
  \item For New Jersey (12 seats), $k = 6$ districts of 2 seats each, or
        $k = 4$ districts of 3 seats each.
  \item For Illinois (17 seats), $k = 5$ districts of 3 seats each
        (with 2 districts of 4 seats for exact seat coverage, or using a
        mixed-magnitude design).
\end{itemize}

We implement two seat-allocation rules within each MMD: Single Transferable
Vote (STV, also known as Proportional Ranked Choice) and Open-List
Proportional Representation (OLPR, the system used in Sweden, Finland, and
Brazil).  These two systems represent the major MMD design choices available:
STV is voter-centric (voters rank individual candidates), while OLPR is
party-centric (voters choose a party list with an optional candidate preference).

\subsection{Why Algorithmic MMD Generation Matters}

Computational redistricting tools have focused almost entirely on
single-member districts.  The algorithmic MMD question has been left to
political scientists, who typically draw illustrative examples by hand
\citep{rush2000redistricting}.  Our contribution is to place MMD generation
on the same principled, reproducible algorithmic footing as single-member
redistricting.  This enables:

\begin{enumerate}
  \item \emph{Systematic comparison}: MMD and single-member plans generated by
    the same algorithm from the same data, making the district-magnitude
    effect cleanly separable from boundary-drawing effects.
  \item \emph{Legal baselines}: Courts and commissions evaluating MMD reform
    proposals need computational plans to compare against; we provide the
    methodology.
  \item \emph{VRA analysis}: The Voting Rights Act imposes specific
    requirements on minority representation; MMDs interact with these
    requirements in complex ways that require systematic analysis.
\end{enumerate}

\subsection{Paper Structure}

Section~2 reviews STV and OLPR mechanics and the MMD compactness literature.
Section~3 describes the bisect $k$-adjustment for MMDs and the per-district
balance-tolerance formula.  Section~4 presents results for the three state
case studies.  Section~5 addresses the constitutional and statutory legal
landscape.  Section~6 concludes.
